% chapters/11/section-06.tex
\section{Menutup Blok Python}
\label{sec:menutup-blok-python}

Bab ini menutup \CPMK{2} sekaligus menutup seluruh Blok~I dan Blok~II, yaitu
seluruh bagian buku yang dikerjakan dengan Python. Bagian ini menyimpulkan apa
yang telah dicapai dan apa yang menyusul.

\subsection{Mini Proyek yang Dihasilkan}

Mini proyek AI pada RPS menuntut tiga keluaran, dan ketiganya telah tersedia
pada bab ini.

\begin{itemize}
  \item \textbf{Kode program}, yaitu mini permainan yang memuat NPC bermesin
        keadaan yang mencari jalurnya sendiri. Kode itu tersusun atas berkas
        \code{ch09_fsm.py}, \code{ch10_astar.py}, dan
        \code{ch11_mini_proyek.py}.
  \item \textbf{Ringkasan desain}, yaitu diagram keadaan pada
        Gambar~\ref{fig:diagram-keadaan} dan diagram arsitektur pada
        Gambar~\ref{fig:arsitektur-integrasi}.
  \item \textbf{Laporan \istilah{playtest}}, yaitu catatan hasil ketiga
        skenario uji pada Tabel~\ref{tab:tiga-skenario} beserta perbaikan yang
        dirumuskan pada Tabel~\ref{tab:perbaikan}.
\end{itemize}

Penilaian mini proyek itu memakai rubrik pada Lampiran~A, dengan bobot
\CPMK{2} sebesar 10 persen dari nilai akhir mata kuliah. Komponen penilaiannya
meliputi implementasi mesin keadaan, implementasi pencarian jalur, penggabungan
keduanya, pengujian beserta laporannya, dan presentasi.

\subsection{Perkakas yang Telah Lengkap}

Seluruh perkakas Blok~I dan Blok~II telah lengkap. Tabel~\ref{tab:perkakas-python}
memuatnya.

\begin{table}[htbp]
\centering
\caption{Perkakas blok Python yang telah lengkap beserta bab yang membahasnya}
\label{tab:perkakas-python}
{\setlength{\tabcolsep}{4pt}\footnotesize%
\begin{tabularx}{\textwidth}{@{}L{3.6cm} L{2.6cm} Y@{}}
\toprule
\headrow \textbf{Perkakas} & \textbf{Dibahas pada} & \textbf{Kegunaan} \\
\midrule
Lingkungan kerja dan repositori & Bab~III &
Menjalankan program dan menyimpan riwayat pekerjaan. \\
\addlinespace
Siklus permainan dan pengendalian waktu & Bab~IV &
Menjalankan permainan dengan laju yang terkendali. \\
\addlinespace
Vektor, masukan, dan gerak & Bab~V &
Menggerakkan benda dengan kecepatan yang sama pada setiap komputer. \\
\addlinespace
Gravitasi, lompatan, dan pantulan & Bab~VI &
Memberi benda sifat jatuh, melompat, dan memantul. \\
\addlinespace
Kelas \code{Sprite}, wadah \code{Group}, dan tabrakan & Bab~VII &
Mengelola banyak benda dan menangani persentuhannya. \\
\addlinespace
Perilaku mengejar & Bab~VIII &
Membuat benda menentukan geraknya sendiri. \\
\addlinespace
Mesin keadaan & Bab~IX &
Menyusun perilaku yang beragam beserta aturan perpindahannya. \\
\addlinespace
Pencarian jalur A* & Bab~X &
Menemukan rute pada peta berhalangan. \\
\addlinespace
Penggabungan keduanya dan penilaian mutu & Bab~XI &
Menghubungkan niat dengan rute, lalu mengukur dan memperbaiki mutunya. \\
\bottomrule
\end{tabularx}}
\end{table}

\subsection{Apa yang Menyusul}

Capaian yang telah Anda peroleh bukan perkakas yang akan ditinggalkan. Pada
Blok~III dan Blok~IV, pekerjaan berpindah ke Unity, dan seluruh gagasan pada
tabel di atas tetap berlaku.

Yang berubah hanyalah \textbf{letak tanggung jawabnya}. Siklus permainan,
perhitungan fisika, deteksi tabrakan, dan penyajian diserahkan kepada mesin.
Waktu Anda beralih kepada pekerjaan yang tidak dapat dikerjakan mesin, yaitu
menentukan desain level, menyusun pengalaman bermain, dan menilai mutu hasilnya.

Pengetahuan tentang mekanisme yang telah Anda bangun sendiri itulah yang
membuat peralihan itu berjalan mulus. Ketika perilaku permainan pada Unity
terasa janggal, Anda memiliki bekal untuk menelusuri sebabnya: apakah
pengaturan komponennya yang keliru, apakah mesinnya yang bekerja dengan cara
yang tidak Anda duga, atau apakah rancangan levelnya yang memang kurang tepat.

\begin{catatabox}
Bab~XII sampai Bab~XVI membahas Unity, dan Bab~XVII menutup buku ini dengan
refleksi atas perbandingan kedua pendekatan. Perbandingan itu akan lebih
bermakna apabila Anda telah merasakan sendiri kedua cara bekerja. Karena itu,
simpanlah catatan praktikum Anda sejak Bab~III; catatan itulah bahan refleksi
yang paling jujur pada akhir semester.
\end{catatabox}
